\section{Dataset and experimental conditions}
\label{sec:dataset}

\paragraph{Ambient condition.}
All experiments are performed at a nominal ambient temperature of
25\,$^\circ$C. We treat the cell as isothermal at this condition in all model
configurations used for this repository.

\paragraph{Cell type.}
The cells are large-format lithium iron phosphate (LFP) cells with nominal
capacity on the order of \TODO{insert value, e.g., 105\,Ah}.

\paragraph{Available tests.}
The raw data folders include characterisation protocols such as OCV--SOC,
GITT, DCIR, HPPC, and RPT, as well as cycling datasets. Table~\ref{tab:tests}
summarises how each test is used in the modelling workflow.

\begin{table}[t]
  \centering
  \caption{Test types and intended use in the workflow.}
  \label{tab:tests}
  \begin{tabular}{ll}
    \toprule
    Test & Primary use in this repo \\
    \midrule
    OCV--SOC & OCV curve calibration / stoichiometry window \\
    DCIR & Ohmic resistance / validation feature (DCIR) \\
    HPPC & Dynamic resistance/RC behaviour; feature extraction \\
    GITT & Step-level metrics; optional apparent diffusivity \\
    RPT & Capacity and incremental capacity (IC) features \\
    Long-term cycling & \SOH trajectories for fine-tuning/validation \\
    \bottomrule
  \end{tabular}
\end{table}

\paragraph{Data preprocessing.}
The repository includes a unified loader that standardises column names,
reconstructs a monotonic time base from absolute timestamps, and normalises
units (mA$\rightarrow$A, mAh$\rightarrow$Ah) for consistency across exports.

\paragraph{Cell cohort selection.}
Because characterisation completeness varies by cell (e.g., some protocols such
as self-discharge and peak power tests are only available for a subset), we
define a default cohort that maximises the ability to calibrate and validate
the electrochemical model before using it as a synthetic trajectory generator.
In this repository, the default selected cohort is recorded in
\texttt{configs/dataset.yaml} and justified in the generated audit report
\texttt{data/processed/cell\_selection\_report.md}. At the time of writing, this
default cohort comprises cells \texttt{0005--0008}, which have a more complete
characterisation set (including GITT, peak-power, and self-discharge tests) than
the remaining cells.
